\naslov{Enačbe drugega reda}

Najpomembnejša enačba drugega reda je drugi Newtonov zakon.
Malce posplošeno ima obliko
\[
  \ddot{x}_i = F_i(x_1, \ldots, x_n)
\]
za $i = 1, \ldots, n$.
Če vpeljemo $p = \dot{x}$ in $q = x$, dobimo sistem prvega reda
\begin{gather*}
  \dot{q}_i = p_i \\
  \dot{p}_i = F_i(q_1, \ldots, q_n)
\end{gather*}
Sistem lahko še posplošimo.
Naj bosta $F, G : M \subseteq \R^{2n} \to \R^n$ preslikavi iz faznega prostora
$M$.
Zanima nas časovni razvoj sistema, ki je podan z enačbami
\begin{gather*}
  \dot{q} = G(q, p), \\
  \dot{p} = F(q, p).
\end{gather*}

Posebej pomembni so sistemi, za katere obstaja funkcija $H : M \to \R$, za
katero velja
\begin{gather*}
  \frac{\partial H}{\partial q_i} = -F_i, \\
  \frac{\partial H}{\partial p_i} = G_i.
\end{gather*}
Taki funkciji pravimo \pojem{Hamiltonian}.

\begin{izrek}
  Če Hamiltonian obstaja, potem je prvi integral sistema.
\end{izrek}

\begin{proof}
  Naj bo $t \mapsto (q(t), p(t))$ neka rešitev sistema.
  Potem imamo
  \[
	\partial_t H(q, p)
	= \partial_q H \dot{q} + \partial_p H \dot{p}
	=
	\begin{bmatrix}
	  - F & F
	\end{bmatrix}
	\cdot
	\begin{bmatrix}
	  G \\ G
	\end{bmatrix}
	= 0
  \]
\end{proof}

\vprasanje{Kaj je Hamiltonian? Dokaži, da je prvi integral.}

\begin{definicija}
  Naj bo podana funkcija $H : M \subseteq \R^{2n} \to \R$.
  Sistem enačb
  \begin{gather*}
	\dot{q} = \partial_p H \\
	\dot{p} = -\partial_q H
  \end{gather*}
  je \pojem{hamiltonski sistem} s \pojem{hamiltonsko funkcijo} $H$.
\end{definicija}

Da bo sistem $\dot{q} = G(q, p), \dot{p} = F(q, p)$ Hamiltonski, mora obstajati
funkcija $H$, za katero velja $\partial_p H = G$ in $\partial_q H = -F$.
Zapisano v drugačni obliki
\[
  \begin{bmatrix}
	\dot{q} \\ \dot{p}
  \end{bmatrix}
  =
  \begin{bmatrix}
	0 & I \\ -I & 0
  \end{bmatrix}
  \cdot
  \begin{bmatrix}
	\partial_q H \\ \partial_p H
  \end{bmatrix}
  =
  \begin{bmatrix}
	0 & I \\ -I & 0
  \end{bmatrix}
  \grad H,
\]
torej
\[
  \grad H
  =
  \begin{bmatrix}
	-F \\ G
  \end{bmatrix}.
\]
Da bo sistem hamiltonski, mora biti torej polje $[-F, G]^T$ potencialno.

\vprasanje{Pod katerim pogojem je sistem hamiltonski? Dokaži.}